\documentclass[11pt]{article}

\usepackage{amsmath, amssymb, amsthm}
\usepackage[margin=1in]{geometry}

\title{A Matrix-Free Method for Large-Scale Quadratic Programs with Constraints}
\author{Thomas Schmelzer \and Martin Stoll \and Stephen Boyd}
\date{}

\begin{document}

\maketitle

\section{Overview}

We consider quadratic programs of the form
\[
\min_x \; \tfrac{1}{2} x^\top P x + c^\top x \quad \text{s.t.} \quad A x \le b,\;\; E x = d,
\]
where \(x \in \mathbb{R}^n\). We are interested in the regime in which \(n\) is large and the matrix \(P \in \mathbb{R}^{n \times n}\) is
not formed explicitly, but is available through matrix-vector products.

Typical examples include problems where
\[
P = X^\top X + \lambda I,
\]
with \(X \in \mathbb{R}^{m \times n}\), or more generally where \(P\) is a linear operator that can be applied efficiently.

\section{Matrix-Free Setting}

We assume access to the following operations:
\begin{itemize}
\item \(v \mapsto P v\),
\item \(v \mapsto A v\), \(v \mapsto A^\top v\),
\item \(v \mapsto E v\), \(v \mapsto E^\top v\).
\end{itemize}

No explicit formation of \(P\), \(X^\top X\), or any dense KKT matrix is required. The cost of applying \(P\) is assumed to be comparable to applying \(X\) and \(X^\top\), for example:
\[
P v = X^\top (X v) + \lambda v.
\]

\section{Optimality Conditions}

The KKT conditions are
\[
\begin{aligned}
P x + c + A^\top \lambda + E^\top \nu &= 0, \\
A x - b &\le 0,\quad \lambda \ge 0,\quad \lambda_i (A x - b)_i = 0, \\
E x - d &= 0.
\end{aligned}
\]

Given an active set \(\mathcal{A}\), we solve the equality-constrained subproblem:
\[
A_\mathcal{A} x = b_\mathcal{A}, \quad E x = d.
\]

The corresponding KKT system is
\[
\begin{bmatrix}
P & A_\mathcal{A}^\top & E^\top \\
A_\mathcal{A} & 0 & 0 \\
E & 0 & 0
\end{bmatrix}
\begin{bmatrix}
x \\ \lambda_\mathcal{A} \\ \nu
\end{bmatrix}
=
\begin{bmatrix}
- c \\ b_\mathcal{A} \\ d
\end{bmatrix}.
\]

\section{Krylov Subspace Solution}

\subsection{MINRES}

The KKT matrix is symmetric but indefinite: the off-diagonal blocks
$[A_\mathcal{A}^\top,\, E^\top]$ introduce negative eigenvalues, so the standard
Conjugate Gradient (CG) method, which requires symmetric positive definiteness, cannot
be applied directly to the full system.  MINRES handles symmetric indefinite systems via
short three-term recurrences and requires only matrix-vector products with the KKT
operator:
\[
(x, \lambda, \nu) \mapsto
\begin{bmatrix}
P x + A_\mathcal{A}^\top \lambda + E^\top \nu \\
A_\mathcal{A} x \\
E x
\end{bmatrix}.
\]
No matrix factorisation is required. Warm starts are used between active-set iterations.

\subsection{Conjugate Gradient in the Null Space}

An alternative that recovers a symmetric positive definite (SPD) system is to eliminate
the equality constraints explicitly and apply CG to the resulting reduced problem.  We
describe the procedure in detail.

\paragraph{Setup}
Collect all currently active equality constraints (active-set inequalities and hard
equalities) into the single system
\begin{equation}\label{eq:combined}
B x = f, \qquad
B = \begin{bmatrix} A_\mathcal{A} \\ E \end{bmatrix} \in \mathbb{R}^{s \times n},\quad
f = \begin{bmatrix} b_\mathcal{A} \\ d \end{bmatrix} \in \mathbb{R}^{s},
\end{equation}
where $s = |\mathcal{A}| + \operatorname{rows}(E)$ is the total number of active
constraints.  We assume $B$ has full row rank ($s < n$) and $P \succ 0$.

\paragraph{Null-space basis via QR}
Compute the full QR factorisation of $B^\top \in \mathbb{R}^{n \times s}$:
\begin{equation}\label{eq:qr}
B^\top = Q \begin{bmatrix} R \\ 0 \end{bmatrix},
\qquad Q = \bigl[\,Q_1 \;\big|\; Q_2\,\bigr],
\end{equation}
where $Q \in \mathbb{R}^{n \times n}$ is orthogonal, $R \in \mathbb{R}^{s \times s}$ is
upper triangular, $Q_1 \in \mathbb{R}^{n \times s}$ spans the column space of $B^\top$,
and $Q_2 \in \mathbb{R}^{n \times (n-s)}$ spans the orthogonal complement.
Because $B Q_2 = (Q_2^\top B^\top)^\top = 0$, the columns of $Q_2$ form an orthonormal
basis for $\ker(B)$.

\paragraph{Particular solution}
A particular solution $x_p$ satisfying $B x_p = f$ is obtained from the QR factors by
triangular back-substitution:
\[
x_p = Q_1\, R^{-\top} f.
\]

\paragraph{Constraint elimination}
Every vector satisfying $Bx = f$ can be decomposed as
\begin{equation}\label{eq:decomp}
x = x_p + Q_2\,\tilde{x}, \qquad \tilde{x} \in \mathbb{R}^{n-s},
\end{equation}
since $B(x_p + Q_2\tilde{x}) = f + B Q_2\tilde{x} = f$.  Substituting
\eqref{eq:decomp} into the quadratic objective and differentiating with respect to
the free variable $\tilde{x}$:
\[
\frac{\partial}{\partial\tilde{x}}\Bigl[
  \tfrac{1}{2}(x_p + Q_2\tilde{x})^\top P(x_p + Q_2\tilde{x})
  + c^\top(x_p + Q_2\tilde{x})
\Bigr] = 0
\]
yields the $(n-s) \times (n-s)$ reduced linear system
\begin{equation}\label{eq:reduced}
\underbrace{Q_2^\top P Q_2}_{\hat P}\,\tilde{x}
= -Q_2^\top(P x_p + c).
\end{equation}
The reduced matrix $\hat{P} = Q_2^\top P Q_2$ is SPD: $P \succ 0$ implies
$\tilde{x}^\top \hat{P}\tilde{x} = (Q_2\tilde{x})^\top P(Q_2\tilde{x}) > 0$ for all
$\tilde{x} \ne 0$ because $Q_2$ has full column rank.  CG is therefore directly
applicable to \eqref{eq:reduced}.

\paragraph{Matrix-free reduced operator}
Each application of $\hat{P}$ to a vector $\tilde{v} \in \mathbb{R}^{n-s}$ proceeds
without forming $\hat{P}$ explicitly:
\begin{enumerate}
  \item $u \leftarrow Q_2\,\tilde{v}$ \hfill (lift to $\mathbb{R}^n$; cost $O(n(n-s))$)
  \item $w \leftarrow Pu = X^\top(Xu) + \lambda u$ \hfill (apply $P$ matrix-free; cost $O(Tn)$)
  \item $\hat{w} \leftarrow Q_2^\top w$ \hfill (project back; cost $O(n(n-s))$)
\end{enumerate}
The dominant per-iteration cost is therefore $O(Tn + n^2)$.

\paragraph{QR setup cost}
Computing the full orthogonal factor $Q$ by applying $s$ Householder reflectors to the
$n \times n$ identity costs $O(n^2 s)$ per active-set step.  When many inequality
constraints are active---e.g.\ long-only with $s = O(n)$---this reaches $O(n^3)$ and
dominates the per-iteration cost, making constraint-eliminated CG slower than MINRES in
that regime.  MINRES, by contrast, incurs no such setup cost and has per-iteration cost
$O(Tn + sn) = O(Tn)$ for $s \ll T$.

\paragraph{Special case: single budget constraint}
When $B = \mathbf{1}^\top$ ($s = 1$), the null-space basis $Q_2$ is the Householder
complement of $\hat{e}_1 = \mathbf{1}/\|\mathbf{1}\|$.  The maps $Q_2\,\tilde{v}$ and
$Q_2^\top w$ can be applied implicitly via the rank-one reflector
$I - 2\hat{e}_1\hat{e}_1^\top$ in $O(n)$, reducing the per-iteration cost from
$O(n^2)$ to $O(Tn)$ and the QR setup to $O(n)$.  This is the approach used in the
minimum-variance companion paper for the single budget constraint.

\paragraph{Recovering the primal solution}
After CG converges to $\tilde{x}^*$, the full weight vector is recovered via
\eqref{eq:decomp}:
\[
x^* = x_p + Q_2\,\tilde{x}^*.
\]

\section{Active-Set Method}

We use a standard active-set method:
\begin{enumerate}
\item Initialize with a feasible point and active set \(\mathcal{A}\).
\item Solve the equality-constrained subproblem.
\item If a constraint is violated, add it to \(\mathcal{A}\).
\item If a multiplier is negative, remove the corresponding constraint.
\item Repeat until optimality conditions are satisfied.
\end{enumerate}

\section{Preconditioning}

Performance depends critically on preconditioning. Possible choices include:
\begin{itemize}
\item Diagonal or block-diagonal approximations,
\item Ridge regularization \(P + \gamma I\),
\item Low-rank approximations \(P \approx B F B^\top + D\).
\end{itemize}

The preconditioner must be inexpensive and matrix-free.

\section{Computational Cost}

Let \(T_P\) denote the cost of computing \(P v\), and \(T_A\) the cost of applying \(A\).

Each Krylov iteration costs
\[
O(T_P + T_A).
\]

The total cost is
\[
O(N_{\text{outer}} \cdot N_{\text{Krylov}} \cdot (T_P + T_A)).
\]

For \(P = X^\top X\), we have \(T_P = O(\text{nnz}(X))\). No cubic operations are required.

\section{Numerical Considerations}

\begin{itemize}
\item Conditioning of \(P\) and the KKT system affects convergence.
\item Ridge regularization improves stability.
\item Scaling of constraints is important.
\item Moderate accuracy is usually sufficient.
\end{itemize}

\section{Comparison with Existing Solvers}

\paragraph{Clarabel}
Clarabel is a modern interior-point method designed for conic optimization. It solves a sequence of KKT systems using sparse direct factorizations. It is robust and highly accurate, but memory and runtime scale with explicit matrix formation and factorization. In large-scale matrix-free settings, this becomes limiting.

\paragraph{OSQP}
OSQP is an ADMM-based operator splitting solver. It factorizes a fixed KKT matrix once and reuses it across iterations, making it efficient for medium-scale problems and repeated solves. However, convergence is linear and performance degrades for ill-conditioned problems or high accuracy requirements. It also requires explicit matrix representations.

\section{Discussion}

The method avoids forming dense matrices and relies only on matrix-vector products. It scales well to large problems but depends on good preconditioning.

Compared to factor models, it does not require low-rank structure, but may be less efficient when such structure is strong and exploitable.

\section{Conclusion}

A matrix-free active-set method combined with Krylov solvers provides a practical approach for large-scale constrained quadratic programs.

\end{document}
